\documentclass[11pt, letterpaper]{article}

\usepackage{amsmath}
\usepackage{enumitem}
\usepackage[margin=1in]{geometry}

\title{Parallelization of Diophantine Solvers}
\author{Hamza Iseric, Elijah Kin, Cameron Yeagle}
\date{March 30, 2026}

\begin{document}

\maketitle

\section*{Background}
Diophantine equations have fascinated since antiquity; Pythagorean triples such as $3^2 + 4^2 = 5^2$ were known as early as 1800 BC. However, it was not until 2002 when the first nontrivial integer solutions to
\begin{equation*}
    a_1^6 + a_2^6 = b_1^6 + b_2^6 + b_3^6 + b_4^6 + b_5^6
\end{equation*}
were found \cite{RM03}. The method relies on the fact that sixth powers are necessarily congruent to either $0$ or $1$ mod $7$. Therefore, the left-hand side is congruent to either $0$, $1$, or $2$ mod $7$, implying that at least three of the right-hand side terms are divisible by $7$. This clever observation allows for a dramatic shrinking of the search space based on solving the equation mod $7^6$.

\section*{Proposal}
We notice that there is potential for parallelization in the algorithm (namely that each iteration of the $a_1$ loop is independent) and propose implementations to exploit this:
\begin{enumerate}[label=\roman*]
    \item Naive C++
    \item C++ with OpenMP
    % I would personally prefer not to use MPI but it's certainly an option
    % \item C++ with MPI
    \item CUDA
\end{enumerate}
We can then compare each implementation, hopefully observing nice parallelization speedup.

\section*{Further Work}
Beyond simply recreating the results of \cite{RM03} and comparing the different parallelization methods, modern hardware should enable us to extend to a larger search space beyond $a_1 \leq 30400$. In addition, it would be interesting to try to adapt these methods to other Diophantine equations. There are still no known solutions to
\begin{equation*}
    a_1^6 = b_1^6 + b_2^b + b_3^6 + b_4^6 + b_5^6 + b_6^6.
\end{equation*}

\bibliographystyle{alpha}
\bibliography{diophantine}

\end{document}
